\chapter{Stroom- en spanningswetten in schakelingen}\label{ch:g8:circuit-laws}

Twee \omterm{def:g2:measuring-length:units}{meters}, twee grootheden --- en nu, eindelijk, de grondwet. Elke
schakeling die je ooit zult bouwen, van een zaklamp tot een stadsnet,
gehoorzaamt vier korte wetten: twee over de boekhouding van de stroom,
twee over die van de \omterm{def:g8:voltage-voltmeter:voltage}{spanning}. Met die wetten en een beetje
algebra kun je aflezingen uitrekenen die je nog niet hebt gemeten ---
en elke meting betrappen die liegt.

\section{De stroomwetten}

\begin{proposition}[Stroomwet voor een serielus]\label{prop:g8:circuit-laws:seriesI}
Langs één lus is de \omterm{def:g8:intensity-ammeter:intensity}{stroomsterkte} overal gelijk:
\[
I_{\text{batterij}} = I_{\text{lamp 1}} = I_{\text{lamp 2}} =
\dots
\]
Eén weg, één mars --- de tolhuisjes van het vorige hoofdstuk, nu een
wet van de grondwet.
\end{proposition}

\begin{proposition}[De knooppuntwet]\label{prop:g8:circuit-laws:junction}
In elk \omterm{def:g5:series-parallel-circuits:parallel}{knooppunt} zijn de binnenstromende \omterm{def:g8:intensity-ammeter:intensity}{stroomsterktes} samen
gelijk aan de uitstromende \omterm{def:g8:intensity-ammeter:intensity}{stroomsterktes} samen: voor een hoofdweg die
zich in twee \omterm{def:g5:series-parallel-circuits:parallel}{takken} splitst,
\[
I = I_1 + I_2 .
\]
Lading wordt bij een splitsing niet geschapen, niet vernietigd en niet
opgeslagen: wat aankomt, vertrekt. (Waar voor elk aantal
\omterm{def:g5:series-parallel-circuits:parallel}{takken} --- tel de aankomsten op, tel de vertrekken op, en de
boeken kloppen.)
\end{proposition}

\begin{omfigure}
\begin{tikzpicture}[scale=0.85]
  \draw (0,0) to[battery1] (0,3.2)
    to[ammeter, l={$I = \qty{0.50}{A}$}] (3.0,3.2);
  \draw (3.0,3.2) to[short] (3.0,2.3)
    to[lamp, l={$I_1 = \qty{0.30}{A}$}] (6.2,2.3)
    to[short] (6.2,3.2);
  \draw (3.0,3.2) to[short] (3.0,4.1)
    to[lamp, l={$I_2 = \qty{0.20}{A}$}] (6.2,4.1)
    to[short] (6.2,3.2);
  \draw (6.2,3.2) to[short] (7.8,3.2) to[short] (7.8,0)
    to[short] (0,0);
  \fill (3.0,3.2) circle (2pt);
  \fill (6.2,3.2) circle (2pt);
\end{tikzpicture}

{\small De knooppuntwet met getallen: \qty{0.50}{A} komt bij de
splitsing aan, $0.30 + 0.20$ vertrekken langs de \omterm{def:g5:series-parallel-circuits:parallel}{takken} --- en
\qty{0.50}{A} voegt zich in het verre \omterm{def:g5:series-parallel-circuits:parallel}{knooppunt} weer samen.}
\end{omfigure}

\section{De spanningswetten}

\begin{proposition}[Spanningswet voor een serielus]\label{prop:g8:circuit-laws:seriesU}
Rond een serielus wordt de \omterm{def:g8:voltage-voltmeter:voltage}{spanning} van de
\omterm{def:g3:first-electric-circuit:battery}{batterij} onder de apparaten verdeeld, en de delen tellen precies
op:
\[
U_{\text{batterij}} = U_1 + U_2 + \dots
\]
Het watervalbeeld houdt de boeken bij: de
\omterm{def:g3:first-electric-circuit:battery}{batterij} tilt de mars over haar volle hoogte, en de apparaten
geven die hoogte onderling uit, druppel voor druppel --- terwijl goede
draden er niets van uitgeven.
\end{proposition}

\begin{proposition}[Spanningswet voor parallelle takken]\label{prop:g8:circuit-laws:parallelU}
Apparaten \omterm{def:g5:series-parallel-circuits:parallel}{parallel} overspannen allemaal dezelfde twee
\omterm{def:g5:series-parallel-circuits:parallel}{knooppunten} --- dus bedient één \omterm{def:g8:voltage-voltmeter:voltage}{spanning} ze allemaal:
\[
U_1 = U_2 = U_{\text{tussen het knooppuntenpaar}} .
\]
Elke \omterm{def:g5:series-parallel-circuits:parallel}{tak} geniet de volle duw tussen de
\omterm{def:g5:series-parallel-circuits:parallel}{knooppunten} --- daarom stralen parallelle lampen op hun volle
nominale helderheid, en daarom biedt elk stopcontact in je huis dezelfde
\qty{230}{V}.
\end{proposition}

\begin{omfigure}
\begin{tikzpicture}[scale=0.85]
  \draw (0,0) to[battery1, l=\qty{4.5}{V}] (0,3.0)
    to[short] (1.8,3.0)
    to[lamp, l={$U_1 = \qty{2.1}{V}$}] (4.6,3.0)
    to[lamp, l={$U_2 = \qty{2.4}{V}$}] (7.4,3.0)
    to[short] (8.8,3.0)
    to[short] (8.8,0)
    to[short] (0,0);
\end{tikzpicture}

{\small De seriespanningswet met getallen: de \qty{4.5}{V} van de
\omterm{def:g3:first-electric-circuit:battery}{batterij} wordt in twee ongelijke delen uitgegeven --- $2.1 + 2.4 =
4.5$, tot op de decimaal.}
\end{omfigure}

\begin{example}[Waarom ongelijke delen?]\label{ex:g8:circuit-laws:unequal}
Twee verschillende lampen \omterm{def:g4:series-circuits:series}{in serie} voeren dezelfde stroom (wet
één) en verdelen de \omterm{def:g8:voltage-voltmeter:voltage}{spanning} toch ongelijk --- \qty{2.1}{V}
tegen \qty{2.4}{V}. Het grootste deel gaat naar de lamp die zich het
meest tegen de mars verzet: de zwaarste werker neemt het grootste
verval. Wat ``zich meer verzetten'' precies betekent, en hoe je het
meet, is het hele onderwerp van het volgende hoofdstuk --- de
wetten van dit hoofdstuk vormen er het toneel voor.
\end{example}

\section{Rekenen voordat je meet}

\begin{method}[Een schakeling oplossen met de vier wetten]\label{met:g8:circuit-laws:solve}
Gegeven een schakeling met enkele bekende aflezingen:
\begin{enumerate}
  \item teken het schema en markeer elke bekende $I$ en $U$;
  \item schrijf de toepasselijke wetten als vergelijkingen op ---
        knooppuntsommen voor stromen, lussommen voor
        \omterm{def:g8:voltage-voltmeter:voltage}{spanningen}, gelijkheid voor $I$ \omterm{def:g4:series-circuits:series}{in serie} en $U$
        \omterm{def:g5:series-parallel-circuits:parallel}{parallel};
  \item los de onbekenden op (de vergelijkingen uit je algebracursus,
        die hun brood verdienen);
  \item toets op gezond verstand: geen lamp die meer ontvangt dan de
        \omterm{def:g3:first-electric-circuit:battery}{batterij} aanbiedt, geen takstroom die groter is dan die
        van de hoofdweg.
\end{enumerate}
\end{method}

\begin{example}[Uitgewerkt: de gemengde schakeling]\label{ex:g8:circuit-laws:worked}
Een \qty{6}{V}-batterij voedt lamp A \omterm{def:g4:series-circuits:series}{in serie} met een
\omterm{def:g5:series-parallel-circuits:parallel}{parallel} paar (B en C). Gemeten: $I_{\text{batterij}} = \qty{0.5}{A}$,
$I_B = \qty{0.3}{A}$, $U_A = \qty{2.5}{V}$. Bereken de rest.
Knooppuntwet: $I_C = 0.5 - 0.3 = \qty{0.2}{A}$. Stroomwet voor
serie: lamp A voert de volle \qty{0.5}{A}. Spanningswet voor de lus:
het paar ontvangt $U_{\text{paar}} = 6 - 2.5 = \qty{3.5}{V}$;
spanningswet voor \omterm{def:g5:series-parallel-circuits:parallel}{parallel}: $U_B = U_C = \qty{3.5}{V}$. Elke
meteraflezing in de schakeling is nu bekend --- vier wetten, elk één
regel algebra.
\end{example}

\begin{example}[De wetten als leugendetector]\label{ex:g8:circuit-laws:liedetector}
Een practicumverslag beweert: \omterm{def:g3:first-electric-circuit:battery}{batterij} \qty{4.5}{V}; twee lampen
\omterm{def:g4:series-circuits:series}{in serie} op \qty{2.1}{V} en \qty{2.9}{V}. De luswet telt de delen op:
\qty{5.0}{V} uitgegeven uit een beurs van \qty{4.5}{V} --- onmogelijk;
iemand heeft een schaal verkeerd afgelezen. Een ander verslag:
\omterm{def:g5:series-parallel-circuits:parallel}{knooppunt} in \qty{0.40}{A}, \omterm{def:g5:series-parallel-circuits:parallel}{takken} uit \qty{0.25}{A} en
\qty{0.20}{A} --- de splitsing ``schept'' \qty{0.05}{A}: afgewezen.
Vóór de wetten zagen foute aflezingen er even goed uit als goede; nu
controleert de grondwet elke bladzijde.
\end{example}

\begin{remark}[Batterijen onder dezelfde wetten]\label{rem:g8:circuit-laws:batteries}
De regel van achter elkaar is de seriespanningswet in haar oudste
kleren: \omterm{def:g3:first-electric-circuit:battery}{batterijen} \omterm{def:g4:series-circuits:series}{in serie} tellen hun \omterm{def:g8:voltage-voltmeter:voltage}{spanningen} op
($1.5 + 1.5 + 1.5 = \qty{4.5}{V}$ --- het geheim van de platte
\omterm{def:g3:first-electric-circuit:battery}{batterij}), en één omgekeerde cel \emph{trekt af}, want haar duw
gaat op aan het bevechten van de andere. De wetten regeren gevers en
uitgevers gelijkelijk; alleen het teken van hun bijdrage verschilt.
\end{remark}

\section{Oefeningen}

\begin{exercise}[$\star$]\label{exo:g8:circuit-laws:1}
Formuleer de vier wetten --- twee voor de stroom, twee voor de
\omterm{def:g8:voltage-voltmeter:voltage}{spanning} --- elk in één regel.
\end{exercise}

\begin{exercise}[$\star$]\label{exo:g8:circuit-laws:2}
Een \omterm{def:g5:series-parallel-circuits:parallel}{knooppunt} ontvangt \qty{0.8}{A} en stuurt \qty{0.55}{A} de
ene \omterm{def:g5:series-parallel-circuits:parallel}{tak} in. Wat voert de andere \omterm{def:g5:series-parallel-circuits:parallel}{tak}?
\end{exercise}

\begin{exercise}[$\star$]\label{exo:g8:circuit-laws:3}
Een \qty{9}{V}-batterij voedt drie lampen \omterm{def:g4:series-circuits:series}{in serie}; twee daarvan
nemen \qty{2.5}{V} en \qty{3.8}{V}. Het deel van de derde?
\end{exercise}

\begin{exercise}[$\star$]\label{exo:g8:circuit-laws:4}
Twee lampen \omterm{def:g5:series-parallel-circuits:parallel}{parallel} over een \qty{12}{V}-batterij:
de \omterm{def:g8:voltage-voltmeter:voltage}{spanning} over elke lamp? Welk alledaags feit over
stopcontacten in huis is dezelfde wet?
\end{exercise}

\begin{exercise}[$\star$]\label{exo:g8:circuit-laws:5}
Welke enkele extra meting zou in de uitgewerkte gemengde schakeling
overbodig zijn geweest, en waarom? (Kies er een en verantwoord het met
een wet.)
\end{exercise}

\begin{exercise}[$\star$]\label{exo:g8:circuit-laws:6}
Twee verschillende lampen \omterm{def:g4:series-circuits:series}{in serie} voeren dezelfde stroom maar
verdelen de \omterm{def:g8:voltage-voltmeter:voltage}{spanning} $1.8$ tegen \qty{2.7}{V}. Welke lamp verzet
zich het meest tegen de mars? Welk komend hoofdstuk meet dat
verzet?
\end{exercise}

\begin{exercise}[$\star$]\label{exo:g8:circuit-laws:7}
Controle: \omterm{def:g3:first-electric-circuit:battery}{batterij} \qty{6}{V}; lampen \omterm{def:g4:series-circuits:series}{in serie} gemeten op
\qty{2.2}{V}, \qty{1.9}{V}, \qty{1.4}{V}. Kloppen de boeken binnen de
eerlijkheid van een meetinstrument?
\end{exercise}

\begin{exercise}[$\star$]\label{exo:g8:circuit-laws:8}
Controle: hoofdweg \qty{0.9}{A}; drie \omterm{def:g5:series-parallel-circuits:parallel}{takken} \qty{0.4}{A},
\qty{0.3}{A}, \qty{0.3}{A}. Vonnis?
\end{exercise}

\begin{exercise}[$\star\star$]\label{exo:g8:circuit-laws:9}
Een \qty{4.5}{V}-batterij, lamp D \omterm{def:g4:series-circuits:series}{in serie} met een
\omterm{ex:g3:first-electric-circuit:switch}{schakelaar} --- en, \omterm{def:g5:series-parallel-circuits:parallel}{parallel} aan lamp D, toont een
\omterm{def:g8:voltage-voltmeter:voltmeter}{voltmeter} \qty{4.5}{V} over de \emph{open}
\omterm{ex:g3:first-electric-circuit:switch}{schakelaar} en \qty{0}{V} over D. Toon aan dat beide aflezingen
de spanningswet voor de lus gehoorzamen.
\end{exercise}

\begin{exercise}[$\star\star$]\label{exo:g8:circuit-laws:10}
Drie identieke cellen achter elkaar laten een lamp
branden; een vierde, omgekeerd, sluit aan in de rij. Totale
\omterm{def:g8:voltage-voltmeter:voltage}{spanning} voor en na, volgens de seriewet met tekens? Wat doet
de lamp?
\end{exercise}

\begin{exercise}[$\star\star$]\label{exo:g8:circuit-laws:11}
Een schakeling met twee \omterm{def:g5:series-parallel-circuits:parallel}{takken}: \omterm{def:g5:series-parallel-circuits:parallel}{tak} één draagt de
lampen E en F \omterm{def:g4:series-circuits:series}{in serie}; \omterm{def:g5:series-parallel-circuits:parallel}{tak} twee draagt lamp G
alleen. \omterm{def:g3:first-electric-circuit:battery}{Batterij} \qty{6}{V}; $U_E = \qty{2.2}{V}$. Bereken $U_F$
en $U_G$, met bij elke stap de wet erbij.
\end{exercise}

\begin{exercise}[$\star\star\star$]\label{exo:g8:circuit-laws:12}
De volledige controle: een \qty{6}{V}-batterij; lamp P
\omterm{def:g4:series-circuits:series}{in serie} met een \omterm{def:g5:series-parallel-circuits:parallel}{knooppunt}; de \omterm{def:g5:series-parallel-circuits:parallel}{takken} Q
en R komen weer samen en keren terug. Gemeten: $I_P = \qty{0.6}{A}$,
$I_Q = \qty{0.45}{A}$, $U_Q = \qty{2.3}{V}$. Bereken elke resterende
stroom en \omterm{def:g8:voltage-voltmeter:voltage}{spanning} in de schakeling ($I_R$, $U_P$, $U_R$), met
bij elke regel de wet erachter --- verzin daarna één extra ``meting''
die een slordige leerling zou kunnen rapporteren en die jouw voltooide
oplossing meteen zou veroordelen.
\end{exercise}

\section{Probleem: de verzegelde mysteriedozen}

\begin{problem}[{Weekendprobleem --- de verzegelde-dozenwedstrijd van
de elektronicaclub; vier wetten tegen vier mysteries; de grote
controle}]\label{pb:g8:circuit-laws:1}
De club verzegelt eenvoudige schakelingen in dozen en laat alleen
meetbussen over. Teams moeten het binnenwerk afleiden --- grondwet in
de hand. De \omterm{def:g3:first-electric-circuit:battery}{batterij} op de tafel is overal \qty{6}{V}.

\medskip\noindent\textbf{Deel I --- Doos A: twee klemmen, twee
lampen.}
Het etiket geeft toe: twee identieke lampen en verder niets, ofwel
\omterm{def:g4:series-circuits:series}{in serie} ofwel \omterm{def:g5:series-parallel-circuits:parallel}{parallel}.
\begin{enumerate}
  \item Aan de \omterm{def:g3:first-electric-circuit:battery}{batterij} gehangen trekt doos A \qty{0.2}{A}; van
        één zo'n lamp alleen op \qty{6}{V} is bekend dat hij ongeveer
        \qty{0.4}{A} trekt. Serie of \omterm{def:g5:series-parallel-circuits:parallel}{parallel}? Betoog met een
        stroomwet.
  \item Voorspel wat de doos in de \emph{andere} schakeling zou
        trekken.
  \item Welke \omterm{def:g8:voltage-voltmeter:voltage}{spanning} ontvangt elke binnenlamp in de werkelijke
        schakeling, volgens welke spanningswet?
  \item In welke schakeling zouden de verborgen lampen op hun volle
        nominale helderheid stralen --- en waarom doen ze dat door het
        kijkgaatje van de doos zichtbaar niet?
\end{enumerate}

\medskip\noindent\textbf{Deel II --- Doos B: de drielampenpuzzel.}
Doos B geeft toe: één lamp (X) \omterm{def:g4:series-circuits:series}{in serie} met een \omterm{def:g5:series-parallel-circuits:parallel}{parallel} paar
(Y, Z --- niet noodzakelijk identiek).
\begin{enumerate}[resume]
  \item Het team meet $I_{\text{batterij}} = \qty{0.5}{A}$ en, bij
        de eigen meetbus van Y, $I_Y = \qty{0.35}{A}$. Vind $I_Z$ en
        $I_X$, met de wetten erbij.
  \item Een \omterm{def:g8:voltage-voltmeter:voltmeter}{voltmeter} over X: \qty{2.8}{V}. Vind de
        \omterm{def:g8:voltage-voltmeter:voltage}{spanning} van het parallelle paar --- en elk van $U_Y$,
        $U_Z$.
  \item Y en Z voeren bij dezelfde \omterm{def:g8:voltage-voltmeter:voltage}{spanning} verschillende
        stromen. Wat onthult dat over de twee lampen, in de taal van
        \cref{ex:g8:circuit-laws:unequal}?
  \item Het team draait Y door een luikje los. Voorspel in welke
        richting $I_{\text{batterij}}$ verandert (stijgen, dalen of
        gelijk blijven), redenerend vanuit wegen en hun verzet.
\end{enumerate}

\medskip\noindent\textbf{Deel III --- Doos C: het bedrog.}
\begin{enumerate}[resume]
  \item De bouwer van doos C rapporteert: ``\omterm{def:g3:first-electric-circuit:battery}{batterij}
        \qty{6}{V}; twee lampen \omterm{def:g4:series-circuits:series}{in serie} erin op \qty{2.5}{V} en
        \qty{4.1}{V}; \omterm{def:g5:series-parallel-circuits:parallel}{knooppunt} erin dat \qty{0.3}{A} ontvangt
        en \qty{0.2}{A} plus \qty{0.15}{A} verstuurt.'' Veroordeel het
        verslag twee keer, één wet per veroordeling.
  \item Van een van de twee gerapporteerde
        spanningsaflezingen blijkt later dat ze eerlijk is. Als
        dat de \qty{2.5}{V} is, wat moet haar partner dan werkelijk
        zijn?
  \item De bouwer bekent dat de knooppuntgetallen ``royaal afgerond''
        waren. Als de twee takaflezingen \qty{0.2}{A} en
        \qty{0.15}{A} eerlijk zijn, wat moet de hoofdweg dan
        werkelijk voeren?
  \item Waarom maken de vier wetten van bedrog met verzegelde dozen
        een verloren spel? Eén zin.
\end{enumerate}

\medskip\noindent\textbf{Deel IV --- De grondwet.}
\begin{enumerate}[resume]
  \item Schrijf het handvest van de club: de vier wetten in vier
        regels, elk met het alledaagse \omterm{prop:g5:mirrors-reflection:image}{beeld} dat haar draagt
        (wegen, tolhuisjes, watervallen, splitsingen).
\end{enumerate}
\end{problem}
